\lecture{14}{2026-11-03}{Fast Polynomial Multiplication via the Number Theoretic Transform}

In Ring-LWE, we work in the polynomial quotient ring
\[
    R_q = \mathbb{Z}_q[x] / (x^d + 1), \quad \text{where } d = 2^k.
\]
Cryptographic operations require computing noisy products \(a(x)s(x) + e(x) \in R_q\).
Naive polynomial multiplication modulo \(x^d + 1\) requires \(O(d^2 \log^2 q)\) bit operations.
To accelerate this to \(O(d \log d \log^2 q)\), we use the Number Theoretic Transform (NTT).

Choose the prime modulus \(q\) such that \(2d \mid (q - 1)\), i.e., \(q - 1 = 2^{k+1} f\) for some integer \(f \ge 1\).
Since \(\mathbb{Z}_q^\times\) is cyclic of order \(q - 1\), there exists a primitive \(2^{k+1}\)-th root of unity \(r \in \mathbb{Z}_q\).
Then \(r^{2^k} = r^d \equiv -1 \pmod q\), so the odd powers \(r, r^3, \dots, r^{2^{k+1}-1}\) are distinct roots of \(x^d + 1\), yielding
\[
    x^d + 1 = \prod_{i=0}^{d-1} \left( x - r^{2i+1} \right) \in \mathbb{Z}_q[x].
\]
By the Chinese Remainder Theorem, evaluation at each root gives the ring isomorphism
\begin{align*}
    R_q = \mathbb{Z}_q[x] / (x^{2^k} + 1)
    &\cong \prod_{i=0}^{2^k-1} \mathbb{Z}_q[x] / (x - r^{2i+1}) \cong \mathbb{Z}_q^{2^k}, \\
    a(x) &\xmapsto{\quad} \left( a(r), a(r^3), \dots, a(r^{2^{k+1}-1}) \right).
\end{align*}
We define the forward transform as this vector of \(2^k\) evaluations:
\[
    \mathsf{NTT}(a, r, k) = \left( a(r), a(r^3), \dots, a(r^{2^{k+1}-1}) \right) \in \mathbb{Z}_q^{2^k}.
\]

To compute this efficiently, write \(a(x) = \sum_{j=0}^{2^k - 1} a_j x^j = a_E(x^2) + x a_O(x^2)\), where
\[
    a_E(y) = \sum_{j=0}^{2^{k-1} - 1} a_{2j} y^j, \qquad a_O(y) = \sum_{j=0}^{2^{k-1} - 1} a_{2j+1} y^j.
\]
Evaluating at an odd power \(r^{2t+1}\) gives
\[
    a(r^{2t+1}) = a_E\left( (r^2)^{2t+1} \right) + r^{2t+1} a_O\left( (r^2)^{2t+1} \right).
\]
Since \(r^2\) has order \(2^k\), the evaluations of \(a_E\) and \(a_O\) at odd powers of \(r^2\) correspond exactly to the half-size transforms \(\hat{a}_E = \mathsf{NTT}(a_E, r^2, k-1)\) and \(\hat{a}_O = \mathsf{NTT}(a_O, r^2, k-1)\).
Using the indexing \(\hat{a}[t] = a(r^{2t+1})\), \(\hat{a}_E[t] = a_E((r^2)^{2t+1})\), and \(\hat{a}_O[t] = a_O((r^2)^{2t+1})\), the periodicity \((r^2)^{2(t + 2^{k-1}) + 1} = (r^2)^{2t+1}\) and the sign flip \(r^{2(t + 2^{k-1}) + 1} = -r^{2t+1}\) yield the butterfly relations
\begin{align*}
    \hat{a}[t]           &= \hat{a}_E[t] + r^{2t+1} \hat{a}_O[t] \pmod q, \\
    \hat{a}[t + 2^{k-1}] &= \hat{a}_E[t] - r^{2t+1} \hat{a}_O[t] \pmod q,
\end{align*}
for each \(t \in \{0, 1, \dots, 2^{k-1} - 1\}\).

This recursive divide-and-conquer strategy is formalized in \Cref{alg:ntt}.

\begin{algorithm}[H]
    \caption{Fast Number Theoretic Transform (Forward NTT)}\label{alg:ntt}
    \begin{algorithmic}[1]
        \Require \(a(x) = \sum_{j=0}^{2^k - 1} a_j x^j \in \mathbb{Z}_q[x] / (x^{2^k} + 1)\), primitive \(2^{k+1}\)-th root of unity \(r \in \mathbb{Z}_q\), integer \(k \ge 0\)
        \Ensure Vector \(\hat{a} = \left( a(r), a(r^3), \dots, a(r^{2^{k+1}-1}) \right) \in \mathbb{Z}_q^{2^k}\)
        \If{\(k = 0\)}
            \State \Return \(\left( a_0 \right)\)
        \EndIf
        \State \(a_E(y) \gets \sum_{j=0}^{2^{k-1} - 1} a_{2j} y^j\)
        \State \(a_O(y) \gets \sum_{j=0}^{2^{k-1} - 1} a_{2j+1} y^j\)
        \State \(\hat{a}_E \gets \mathsf{NTT}(a_E, r^2, k - 1)\)
        \State \(\hat{a}_O \gets \mathsf{NTT}(a_O, r^2, k - 1)\)
        \For{\(t \gets 0\) \textbf{to} \(2^{k-1} - 1\)}
            \State \(\hat{a}[t] \gets \hat{a}_E[t] + r^{2t+1} \hat{a}_O[t] \pmod q\)
            \State \(\hat{a}[t + 2^{k-1}] \gets \hat{a}_E[t] - r^{2t+1} \hat{a}_O[t] \pmod q\)
        \EndFor
        \State \Return \(\hat{a}\)
    \end{algorithmic}
\end{algorithm}

Let \(M_k\) be the number of ring operations in \(\mathbb{Z}_q\) performed by \Cref{alg:ntt} on a polynomial of degree \(2^k - 1\).
The recurrence is
\[
    M_k = 2M_{k-1} + 2 \cdot 2^{k-1}, \quad M_0 = 0,
\]
which solves to \(M_k = k 2^k = d \log_2 d\).
Therefore, the forward NTT runs in \(O(d \log d)\) operations in \(\mathbb{Z}_q\), or \(O(d \log d \log^2 q)\) bit operations.

\section*{Ring-LWE Encryption in the NTT Domain}

Since the NTT is a ring isomorphism, polynomial multiplication in \(R_q\) corresponds to pointwise (Hadamard) multiplication \(\circ\) in the NTT domain \(\hat{R}_q \cong \mathbb{Z}_q^d\).
This yields the accelerated Ring-LWE scheme:
\begin{itemize}
    \item \(\mathsf{KeyGen}()\): Sample \(\hat{a} \xleftarrow{\$} \hat{R}_q\) directly in the NTT domain, and small polynomials \(s, e \leftarrow \chi^d\).
    Compute \(\hat{s} = \mathsf{NTT}(s)\), \(\hat{e} = \mathsf{NTT}(e)\), and
    \[
        \hat{b} = \hat{a} \circ \hat{s} + \hat{e} \in \hat{R}_q.
    \]
    Output \(\mathsf{pk} = (\hat{a}, \hat{b})\) and \(\mathsf{sk} = \hat{s}\).
    \item \(\mathsf{Enc}(\mathsf{pk}, m(x))\): Sample small polynomials \(r, e', e'' \leftarrow \chi^d\).
    Compute \(\hat{r} = \mathsf{NTT}(r)\), \(\hat{e}' = \mathsf{NTT}(e')\), \(\hat{e}'' = \mathsf{NTT}(e'')\), and \(\hat{m} = \mathsf{NTT}(m(x) \lfloor q/2 \rfloor)\).
    Compute
    \[
        \hat{c}_1 = \hat{a} \circ \hat{r} + \hat{e}' \in \hat{R}_q, \qquad \hat{c}_2 = \hat{b} \circ \hat{r} + \hat{e}'' + \hat{m} \in \hat{R}_q.
    \]
    Output the ciphertext \(\hat{c} = (\hat{c}_1, \hat{c}_2)\).
    \item \(\mathsf{Dec}(\mathsf{sk}, \hat{c})\): Compute \(\hat{w} = \hat{c}_2 - \hat{c}_1 \circ \hat{s} \in \hat{R}_q\), and invert \(w(x) = \mathsf{NTT}^{-1}(\hat{w}) \in R_q\).
    Decode each bit \(m_i = 0\) if \(|w_i|_q < q/4\), and \(1\) otherwise.
\end{itemize}
All polynomial multiplications now run in \(O(d \log d)\) ring operations instead of \(O(d^2)\).

\section*{Module-LWE}

While Ring-LWE is efficient, scaling security by increasing the degree \(d\) requires choosing a new modulus \(q\) (to satisfy \(2d \mid q - 1\)) and reconfiguring the NTT.
\emph{Module-LWE} solves this by replacing scalar polynomials with vectors and matrices of polynomials over a fixed ring \(R_q = \mathbb{Z}_q[x] / (x^d + 1)\).

\begin{definition}[Module-LWE Problem]
    Let \(R_q = \mathbb{Z}_q[x] / (x^d + 1)\).
    Given a uniformly random matrix \(A \xleftarrow{\$} R_q^{k \times \ell}\), a small secret vector \(s \leftarrow \chi_s^\ell\), and a small error vector \(e \leftarrow \chi_e^k\), the \emph{Module-LWE} problem asks to recover \(s\) from
    \[
        (A, \, b = As + e \in R_q^k).
    \]
\end{definition}

Setting \(k = \ell = 1\) recovers Ring-LWE, while setting \(d = 1\) recovers standard LWE.
By applying the NTT componentwise, all matrix-vector products remain fast.

\begin{example}[Kyber / ML-KEM]
    Kyber fixes \(d = 256\) and \(q = 3329\), and scales security by varying the matrix dimension \(k = \ell \in \{2, 3, 4\}\):
    \begin{itemize}
        \item Kyber-512 (NIST Level 1): \(k = 2\);
        \item Kyber-768 (NIST Level 3): \(k = 3\);
        \item Kyber-1024 (NIST Level 5): \(k = 4\).
    \end{itemize}
    Secrets and errors are sampled from centered binomial distributions.
    This allows a single, fixed NTT implementation to support all three security levels.
\end{example}